\begin{theorem}[G-19: an explicit rectangular discriminant module and its positive energy]
Let $m\geq3$, put $N=m^2$, and work over $\mathbb C$.
Here $V=\mathbb C^N$ is the argument space, forms lie in
$\operatorname{Sym}^m V^*$, and coefficient functions of degree
$D$ lie in $\operatorname{Sym}^D(\operatorname{Sym}^m V)$.
We use $g\cdot F=F\circ g^{-1}$, hence
$(g\cdot h)(F)=h(F\circ g)$.
Let $k\geq5$ satisfy
\[
 k\leq \operatorname{codim}_{\mathbb A^N}
                 \operatorname{Sing}(\operatorname{per}_m).
\]
In particular $k=5$ is allowed for every $m\geq3$, and $k=6$
is allowed for every $m\geq6$, by the established bounds in R-B17.
Define
\[
 D=k(m-1)^{k-1},\qquad \delta=m(m-1)^{k-1},\qquad
 \lambda=(\delta^k).
\]
There is an explicit highest-weight polynomial $h$ of coefficient
degree $D$ and weight $\lambda$, namely the discriminant of the
restriction to the first $k$ variables, with these properties.

Its generated irreducible module $W$ is contained in the ideal of
the determinant orbit closure, while its restriction to the
permanent orbit closure is nonzero and therefore injective on $W$.
With this explicit dual convention, $W$ has positive polynomial
type $S_\lambda V$. For normalized Haar measure, and any
$U(N)$-invariant Hermitian inner product on $W$, one has
\begin{equation}\label{G19:Haar}
 \int_{U(N)}|h(g\cdot F)|^2\,dg
   =\frac{\|h\|^2}{\dim W}\sum_{j=1}^{\dim W}|h_j(F)|^2,
 \qquad
 \dim W=\prod_{i=1}^k\prod_{j=k+1}^{N}
                     \frac{\delta+j-i}{j-i},
\end{equation}
where $(h_j)$ is an orthonormal basis of $W$.
This integral is zero at $F=\det_m$ and strictly positive at
$F=\operatorname{per}_m$.

There is also a positive integer Gaussian certificate. Normalize
the discriminant to have primitive integer coefficients, let
$A:\mathbb C^k\longrightarrow\mathbb C^N$ have independent
matrix entries, and write
\[
 H_m(A)=\operatorname{Disc}_{k,m}(\operatorname{per}_m\circ A)
        =\sum_{|\alpha|=mD}c_\alpha A^\alpha.
\]
Then $H_m$ is a nonzero integer polynomial, and for independent
standard complex Gaussian entries $Z$ normalized by
$\mathbb E|Z_{ij}|^2=1$,
\begin{equation}\label{G19:Gaussian}
 \mathbb E|H_m(Z)|^2=\sum_\alpha c_\alpha^2\alpha!\ \geq1.
\end{equation}
For the determinant in place of the permanent, the polynomial
and expectation are identically zero. There exists an integer
matrix $A$ with all entries in $\{0,1,\ldots,mD\}$ and
$H_m(A)\ne0$.

This separates a specified irreducible copy, not the total
multiplicities of its isomorphism type. It proves neither
$m_\lambda(\operatorname{per}_m)>m_\lambda(\det_m)$ nor
$m_\lambda(\det_m)=0$. Moreover the analogous degree-$n$
discriminant construction vanishes on every positively padded
form $\ell^{n-m}f_m$, so it supplies no such separator when
$n>m$.
\end{theorem}

\paragraph{Known inputs and the precise deduction.}
The geometric determinant input is Alper, Bogart, and Velasco
(2017), Theorem~1.7, author version printed p.~3, with proof
on p.~5: every linear determinant in at least five variables
defines a singular projective hypersurface. The permanent
singular-locus bound is the cited R-B17 Theorem~4.14, strengthened
by its Theorem~4.23 for $m\geq6$. The discriminant and its
invariance are Mukai (2003), Definition~5.20 and Corollary~5.21,
printed p.~170. Its degree and change-of-variable weight follow
from Benoist (2012), Theorem~1.3, equations (1.4)--(1.6),
printed pp.~1190--1191, specialized to one equation; the
coefficient-degree specialization is also Example~1.8, p.~1192.
The second displayed formula of that example omits the factor
$d_1$ (written $m$ here) in the variable-change degree; we use equation~(1.4),
Theorem~1.3(1.6), and scalar dilation for the correct weight.
These known theorems are not reproved. The deduction below
identifies a concrete rectangular module and its Haar/Gaussian
certificates, and states explicitly why this does not compare
the total multiplicities or survive padding. No historical
priority claim is made for the underlying discriminant separator.

\begin{proof}
Let $X=\{\operatorname{per}_m=0\}\subset\mathbb P^{N-1}$.
The codimension hypothesis implies that a general projective
$(k-1)$-plane misses $\operatorname{Sing}X$: its expected
intersection dimension is at most
$(k-1)+(N-1-k)-(N-1)=-1$.
Iterated Bertini on the smooth locus then makes its intersection
with $X$ smooth. This is the same linear-section argument used
in the cited proof of Alper--Bogart--Velasco Theorem~1.7, now
applied to the permanent's singular-locus bound. Consequently
\[
 \operatorname{Disc}_{k,m}(\operatorname{per}_m\circ A)\ne0
\]
for $A$ in a nonempty Zariski-open subset of the space of linear
maps. In contrast, the cited determinant theorem makes
$\operatorname{Disc}_{k,m}(\det_m\circ A)=0$ for every $A$.
If a restriction is the zero form, its discriminant is zero as
well, so no exceptional lower-rank map changes this conclusion.

The discriminant has coefficient degree $D=k(m-1)^{k-1}$.
Under substitution $x\mapsto Bx$ it transforms by
$\det(B)^\delta$, where $k\delta=mD$; thus
$\delta=m(m-1)^{k-1}$. Regard this determinant-character vector
as an element of the plethysm in $N$ variables using the first
$k$ coordinate vectors. It is killed by the raising operators
inside $\mathrm{GL}_k$, because it is a determinant character,
and by the raising operators from later coordinate indices,
because no such index occurs. It is therefore a highest-weight
vector of weight $(\delta^k,0^{N-k})$. Semisimplicity shows that
the module generated by this one vector is one irreducible copy
$W$ of that type, even if the ambient plethysm contains other
copies of the same type. The displayed Weyl dimension formula
is the usual formula with all equal-row factors canceled.

Every translate of $h$ vanishes on $\det_m$, and hence $W$
vanishes on its entire orbit closure. A full-rank map $A$ with
smooth permanent restriction can be completed to an invertible
change of the $N$ variables. Therefore some translate of $h$
is nonzero at $\operatorname{per}_m$. The restriction map from
$W$ to the permanent coordinate ring is equivariant and nonzero;
irreducibility makes it injective.

Apply the Schur-orthogonality projector identity used in R-02
and R-05 to this particular irreducible copy. It gives
\eqref{G19:Haar}. At least one basis vector evaluates nontrivially
at the permanent, so its energy is strictly positive. Every
basis vector evaluates to zero at the determinant. This proof
uses compact integration only after identifying the actual
copy and its two restriction maps.

Choose the primitive integer equation of the discriminant,
as permitted by the rational discriminant construction in
Benoist (2012), Theorem~1.7, p.~1191. The coefficients of
$\operatorname{per}_m\circ A$ are integer polynomials of degree
$m$ in the entries of $A$, so $H_m$ is integral and homogeneous
of degree $mD$. It is nonzero by the smooth-section argument.
The exact complex-Gaussian monomial orthogonality in R-06 gives
\eqref{G19:Gaussian}; some nonzero integer coefficient has
square at least one, and every $\alpha!$ is a positive integer.
Finally a nonzero polynomial of total degree at most $mD$
cannot vanish on the entire grid
$\{0,\ldots,mD\}^{kN}$, by successive univariate interpolation.
This proves the finite integer witness bound. It does not
give a polynomial-time search through that grid or an efficient
evaluation procedure for its permanent coefficients.

For the two limitations, the kernel on one specified copy $W$
does not determine the kernel on the whole $\lambda$-isotypic
space: other copies may survive on the determinant, while other
independent copies may vanish on the permanent. Thus no total
multiplicity inequality follows from this argument.
If $n>m$, every nonzero $k$-variable restriction of
$\ell^{n-m}f_m$ has a linear factor and a positive-degree
remaining factor. If the linear factor repeats, its hypersurface
is singular along that hyperplane. Otherwise the hyperplane
meets the remaining positive-degree hypersurface in projective
space, and all first derivatives of the product vanish at their
intersection. Here $k\geq5$ ensures that this intersection is
nonempty. A zero restriction again has zero discriminant.
Hence the degree-$n$ discriminant vanishes on the entire padded
orbit closure. The displayed module and energy cannot remove
this padding obstruction by themselves.
\end{proof}

\paragraph{An explicit smooth five-variable cubic slice.}
For $m=3$, the following matrix provides a concrete witness to
the permanent nonvanishing in G-19:
\[
 M(x)=\begin{pmatrix}
 x_1+x_2&x_3&x_4\\
 x_4&x_0&x_1\\
 x_2+x_3&x_1&x_0+x_2
 \end{pmatrix}.
\]
Writing $f=\operatorname{per}M$, direct expansion gives
\begin{align*}
 f={}&x_0^2x_1+x_0^2x_2+x_0x_1x_2+x_0x_2^2
        +x_0x_2x_4+2x_0x_3x_4\\
    &+x_1^3+x_1^2x_2+x_1x_2x_3+x_1x_3^2
        +x_1x_4^2+x_2x_3x_4.
\end{align*}
Here is a finite integer certificate of its smoothness, with the
complete matrix and two independently written reconstructions
included in the work archive. Put
$S=\mathbb Z[x_0,x_1,x_2,x_3,x_4]$, and form the coefficient
matrix of
\[
 \bigoplus_{j=0}^4 S_4\longrightarrow S_6,
 \qquad (g_j)_j\longmapsto\sum_{j=0}^4g_j\partial_{x_j}f.
\]
Order each monomial basis by decreasing lexicographic order on
its exponent tuple. The matrix has $210$ rows and $350$ columns;
column $70j+b$ corresponds to the $b$th degree-four monomial
times $\partial_{x_j}f$, with all indices starting at zero.
Keep all rows and select the following multiplier indices in
each derivative block:
\begin{center}
\begin{tabular}{r p{0.82\linewidth}}
\toprule
$j$ & Selected values of $b$ (inclusive intervals)\\
\midrule
0 & $0$--$69$\\
1 & $0$--$12$, $14$--$22$, $24$--$28$, $30$--$31$,
    $34$--$42$, $44$--$48$, $50$--$51$, $54$--$58$,
    $60$--$61$, $64$--$65$, $69$\\
2 & $0$--$12$, $15$--$22$, $25$--$28$, $35$--$42$,
    $45$--$48$, $55$--$58$\\
3 & $0$--$10$, $12$, $15$--$20$, $22$, $25$,
    $35$--$40$, $42$, $45$\\
4 & $0$--$8$, $10$, $12$, $15$--$18$, $35$\\
\bottomrule
\end{tabular}
\end{center}
The resulting $210\times210$ determinant, in the displayed
increasing row and column orders, is $94$ modulo $101$.
It is therefore a nonzero integer. After extension to
$\mathbb C$, every degree-six monomial belongs to the ideal
generated by the five partial derivatives. A common nonzero
projective zero of those derivatives would make all degree-six
monomials vanish, including the sixth power of a nonzero
coordinate, which is impossible. Thus $\{f=0\}\subset\mathbb P^4$
is smooth. This proof uses neither a numerical tolerance nor an
inference from a randomly sampled point.

The author reconstruction uses exact polynomial expansion; the
independent reconstruction multiplies all six permutations in
integer polynomial dictionaries, differentiates and convolves
the dictionaries, reproduces every one of the $73{,}500$ matrix
entries modulo $101$, and separately eliminates the selected
minor. The files are
\texttt{work/root/g19\_cubic/certificate.json},
\texttt{work/root/g19\_cubic/macaulay\_degree6.u8},
\texttt{work/root/check\_g19\_cubic\_slice.py}, and
\texttt{work/isolation/check\_G19\_cubic\_independent.py}.
The row-major byte matrix has SHA-256
\texttt{c739e7d83b6c43a3e48e6f20c92667fd77282af7e1240c3806d55989af0928e5}.

Consequently the specified-copy separator at $m=3$ has
coefficient degree $D=80$ and rectangular weight $(48^5)$.
Five is the minimum possible number of rows for such a
highest-weight separating equation: R-11 shows that the
four-variable determinant-cubic image is dense in the whole
cubic coefficient space. A highest-weight equation with at
most four rows lies in the corresponding coefficient
subalgebra, so vanishing on the determinant forces it to be
zero. G-19 and the displayed smooth slice supply a nonzero
five-row equation. This row-minimality assertion concerns
specified polynomial equations, and does not assert a total
multiplicity gap at weight $(48^5)$.
